% !TEX root = ../../技术文档.tex
\section{仿真环境与验证方法}

仿真验证按运算单元、功能模块和系统接口逐层展开，覆盖CPU、NPU、系统互连及图像采集与显示链路。
测试平台通过接口激励、参考结果比对和时序检查，核对计算结果、请求与响应关系以及异常处理过程。
本节说明共用工具与测试方法，各模块的具体用例和结果分别列于后续附录。

\subsection{工具配置}

仿真环境部署在Windows平台，使用PowerShell和Python组织源码编译、固件生成、仿真运行及结果整理。
ModelSim承担RTL仿真和波形调试，Icarus Verilog用于模块的功能回归。主要工具配置见表\ref{tbl:simulation-tools}。

\begin{reporttable}
\begin{longtable}{@{}p{0.25\linewidth}p{0.24\linewidth}p{0.44\linewidth}@{}}
\caption{仿真工具配置}\label{tbl:simulation-tools}\\
\toprule
工具 & 版本或配置 & 用途 \\
\midrule
\endfirsthead
\caption[]{仿真工具配置（续）}\\
\toprule
工具 & 版本或配置 & 用途 \\
\midrule
\endhead
\bottomrule
\endfoot
ModelSim & SE-64 2020.4 & Verilog/SystemVerilog编译、四态仿真和WLF波形查看 \\
Icarus Verilog & 12.0 devel & 数字逻辑、接口及地址译码的功能回归 \\
RISC-V GCC & 15.2.0，ILP32 ABI & 编译裸机测试固件，生成ELF、存储镜像和反汇编文件 \\
Pango PDS仿真库 & 2022.2-SP6.4配套库 & 提供存储器、乘加单元等器件原语模型 \\
Python & 3.13.13 & 生成测试数据和参考结果，调度仿真并汇总日志 \\
\end{longtable}
\end{reporttable}

表中版本对应CPU及组件验证环境。附录F的人脸检测整机联合仿真使用ModelSim 10.5、RISC-V GCC 15.2.0-1与Python 3.13.5，配套PDS版本为2022.2-SP6.4。
各项验证按其原始运行记录保留工具版本和输入校验值。

各测试顶层使用明确的源码清单、包含目录和参数配置，按依赖关系编译待测模块、测试台及所需原语模型。
CPU用例记录发射宽度、Cache和分支预测配置；NPU用例记录算子参数、并行度及模型数据；Video用例记录图像尺寸、时钟关系、存储布局和显示模式。
测试用例中的内存容量、地址范围和响应延迟随配置保存，用于确定该次运行的适用范围。

定向裸机固件按所需指令集和ILP32 ABI编译，载入测试台的程序存储模型。
指令集回归使用对应测试集的ELF及参考签名；模型计算使用固定的输入张量、权重和量化参数。
运行记录保存源码与输入文件的SHA-256、工具版本、编译命令和仿真参数，使结果能够对应到具体配置。

仿真输出包括编译与运行日志、检查计数、结果文件及WLF/VCD波形。
波形查看日志与批处理运行日志分别保存，附录中的仿真截图来自已保存的波形文件。
周期统计按测试台时钟和指定事件窗口计算，仿真器在计算机上的运行耗时单独记录。

\subsection{测试平台}

测试平台由激励产生、待测模块、外部接口模型和监测比对逻辑组成。
激励端产生时钟、复位、指令或数据请求；接口模型模拟存储和外设的响应；监测逻辑在有效握手及结果有效时检查数据和状态。
跨时钟域测试分别产生核心时钟与像素时钟，并设置不同相位、暂停及复位场景，检查异步FIFO和状态握手过程。

CPU单元测试直接驱动ALU、乘除法器、寄存器堆和CSR接口，将运算结果与独立计算值或预期寄存器状态比较。
模块测试连接取指前端、发射通路、LSU及Cache的请求与响应端口，插入响应延迟和背压，检查保持、旁路、冲刷及错误处理。
整核测试由CPU执行固件，监测指令PC、退休事件、寄存器写回、访存和异常入口；指令出队、发射与退休分别计数，避免将暂停期间保持的有效信号重复计入。
对写回式Cache，测试台分别检查CPU可见数据和下级存储数据，并在相应写回完成后核对下级内容。

NPU计算测试使用与模型包一致的整数参考，逐项核对卷积、辅助算子、重定标和输出张量。
参考计算明确输入符号、累加精度、补边、舍入与饱和规则；存储模型检查访问范围、字节使能、数据排列及重复或遗漏写入。
任务控制测试另外检查启动参数锁存、输出写入完成、完成记录保持和标识匹配的ACK，区分算子结束、结果就绪与完成记录释放。

Video测试由可重复的像素序列产生原图及预期输出，分别核对缩放、通道排列、补边、DMA写入和显示合成。
参考像素根据源坐标与规定的整数运算生成，并随输出流水延迟核对RGB、同步信号、坐标和帧标识。
帧管理测试同时跟踪槽号、帧号及代次，检查显示记录接收、首个有效像素实际输出、旧帧引用解除和释放确认的先后关系。
连续帧场景在采集、计算控制和显示读取重叠时核对数据，检测提前复用、重复写入和显示欠载。

系统接口测试通过CPU固件或总线激励访问MMIO，结合字节存储模型检查CPU、NPU与Video之间的数据传输。
AXI写地址、写数据及写响应通道分别施加背压，读通道独立设置返回延迟；测试包含跨边界访问、错误响应和完成事件延后等场景。
已接纳事务的地址、数据和响应对应关系持续跟踪到事务结束，错误注入用例用于检查故障保持、恢复和错误路径隔离。

每个用例设置有限运行周期和预期结束事件，执行脚本同时读取自检标志、错误信息及退出状态，并保存对应日志。
测试记录注明采用的外部存储模型和时钟配置；RTL功能仿真、布局布线时序分析及板级测量分别记录。
